%! TEX root = thesis.tex

\chapter{Experimental Setup}
\label{sec:experimental_setup}

\section{Models}
\label{sec:models}

We evaluate two state-of-the-art LLMs that represent the current frontier of
large-scale pretraining:

\begin{itemize}
  \item \textbf{DeepSeek~V3}~\cite{deepseek-ai_deepseek-v3_2024}: a 671B
    parameter mixture-of-experts model (37B active parameters per
    token), trained on 14.8T
    tokens using 2048 NVIDIA H800 GPUs, consuming 2.79M GPU-hours.
  \item \textbf{Kimi~K2}~\cite{kimi_team_kimi_2025}: a 1T parameter
    mixture-of-experts model
    (32B active), trained on 15.5T tokens with a minimum 256-GPU model-parallel
    group.
\end{itemize}

These two models differ substantially in scale: DeepSeek~V3 uses 8$\times$ more
GPUs, allowing us to examine how GPU count affects the savings-overhead
trade-off.

\section{Hardware Constants}
\label{sec:hardware-constants}

We use the following hardware parameters for all simulations:

\begin{table}
  \centering
  \caption{Hardware constants used in all simulations.}
  \label{tab:constants}
  \small
  \begin{tabular}{lr}
    \toprule
    Parameter & Value \\
    \midrule
    GPU training power & 700~W (NVIDIA H800) \\
    GPU idle power & 60~W \\
    Power usage effectiveness & 1.27 \\
    Checkpoint pause time & 148.8~s \\
    Checkpoint resume time & 0~s (async I/O) \\
    \bottomrule
  \end{tabular}
\end{table}

GPU training power is based on the NVIDIA H800 thermal design power
thermal design power~\cite{h800_specs}.  GPU idle power is estimated
from published hardware
measurements.  Power usage effectiveness is taken from DeepSeek's DGX
H800 infrastructure
report~\cite{jegham_hungry_2025}.  Checkpoint overhead is based on PyTorch
Distributed synchronous checkpoint measurements for a 7B model, with the
conservative assumption that checkpoint resume is overlapped with asynchronous
I/O~\cite{pytorch_checkpoint}.

\section{Regions and Data}
\label{sec:regions-data}

We evaluate five electricity grid regions: Germany (DE), Sweden (SE), Italy
(IT), the United States (US), and China (CN).  Carbon intensity data is
obtained from Electricity Maps at 5-minute resolution for the year 2025.
Table~\ref{tab:grid} in Chapter~\ref{sec:methodology} summarizes the grid
characteristics of each region.

\section{Optimizer Settings}
\label{sec:optimizer-settings}

All optimization runs use the following configuration:

\begin{table}
  \centering
  \caption{Adaptive optimizer configuration.}
  \label{tab:optimizer-config}
  \small
  \begin{tabular}{lr}
    \toprule
    Parameter & Value \\
    \midrule
    Grid resolution $N$ & 10 points per axis \\
    Maximum iterations $K$ & 10 \\
    Overhead budget $B$ & 200\% \\
    Shrink factor $s$ & 0.45 \\
    Trade-off parameter $\alpha$ & 1.0 (pure savings) \\
    $\theta_p^{\max}$ (Sweden) & 100~gCO$_2$eq/kWh \\
    $\theta_p^{\max}$ (others) & 800~gCO$_2$eq/kWh \\
    Start-date samples $D$ (optimized) & every 7th day \\
    Historical years & 2025 \\
    \bottomrule
  \end{tabular}
\end{table}
